% Exposé XI. What is analogy, and what is open.
% Sources: theory/framework.md (Exposé XI), site/sections/12-predictions.html (second half: "What the borrowed words
% mean", "Open problems", "Debts", "The referee record"), theory/propositions.json (status, round2 and public statements).
% Corrections applied from propositions.json: DoRA's generic dimension (Prop II.8) is now proved, and Obs IV.4 has a full
% proof; the remaining sketch is corollary 4 of Thm V.3 (unconstrained generators), completed in Appendix C.
\section{What is analogy, and what is open}\label{sec:analogy}

This exposé gives a plain account of which borrowed words carry theorems, what is open, whose results are re-derived
here, and what the referees changed. Referees repaired 27 of the \nPropositions{} results, and a second review narrowed
19 of those again.

% ----------------------------------------------------------------------------------------------------------------------
\subsubsection{What the borrowed words mean}\label{sec:analogy-words}

Words borrowed from geometry and physics compress a great deal, and each also carries connotations no theorem here
supports. Table~\ref{tab:analogy} gives, for seven such words, the sense each has in this paper, the results behind it,
and what it leaves out.

\begingroup
\footnotesize
\renewcommand{\arraystretch}{1.18}
\setlength{\LTpre}{6pt}\setlength{\LTpost}{8pt}
\begin{longtable}{@{}>{\raggedright\arraybackslash}p{2.1cm}>{\raggedright\arraybackslash}p{3.75cm}>{\raggedright\arraybackslash}p{4.5cm}>{\raggedright\arraybackslash}p{3.75cm}@{}}
\caption{Analogy versus theorem: seven borrowed words, their meaning here, what is proved under each, and what is not
claimed.}\label{tab:analogy}\\
\toprule
\headrow \textbf{Word} & \textbf{Meaning here} & \textbf{Proved} & \textbf{Not claimed}\\
\midrule
\endfirsthead
\multicolumn{4}{@{}l}{\footnotesize\itshape Table~\ref{tab:analogy}, continued}\\
\toprule
\headrow \textbf{Word} & \textbf{Meaning here} & \textbf{Proved} & \textbf{Not claimed}\\
\midrule
\endhead
\midrule
\multicolumn{4}{@{}r}{\footnotesize\itshape continued on the next page}\\
\endfoot
\bottomrule
\endlastfoot
\textbf{Grothendieck} & The relative point of view: a method is a map $\rho:(Q,q_0)\to(\Theta,\theta_0)$, methods
  form the pointed slice $\Meth(\Theta,\theta_0)$, and changing $\theta_0$ is base change. & Frozen model initial and full
  fine-tuning terminal; products are fibre products in $\Diff$ and can fail to exist (\thmref{Obs.}{I.4}); a symmetric
  monoidal sum (\thmref{Obs.}{I.7}); automorphisms by transport (\thmref{Prop.}{II.13}); pushforward along
  function-preserving maps (\thmref{Def.}{V.1}(a)). & No schemes, sheaves, descent theory or stacks. Quantisation and
  rebasing borrow only the word ``base change'', and idle constructions such as $\int\Meth\to\mathbf{PNet}$ are left
  out. ``Exposé'' and the rising sea are homage.\\
\textbf{Erlangen} & After \citet{xklein1872vergleichende}, a method's geometry is its covariance group $\mathcal G_M$
  (\thmref{Def.}{II.14}, \bgref{D.17}); what a group-type method can never change is a complete invariant of the
  action. & Complete invariants of the orthogonal, torus and Hadamard actions (\thmref{Thm}{II.5}); exactly orthogonal
  input-side methods keep $K_{\mathrm{out}}$ and the spectrum through any number of merges (\thmref{Cor.}{II.6},
  \thmref{Thm}{V.3}); covariance groups of named methods (\thmref{Prop.}{II.9}(e), \thmref{Prop.}{II.16}). & For cones
  such as LoRA's image the word only names the shape and its covariance. Keeping an invariant concerns reach; that it
  causes knowledge retention is not shown.\\
\textbf{Noether} & The pattern that turns a continuous symmetry of $\rho$ into a conserved quantity. &
  \thmref{Thm}{III.5}, for gradient flow driven by any signal $\Lambda D\rho\T\gamma$, a mild extension of
  \citet{zhao2022symmetries}; for LoRA the charges are the entries of $\Phi$. & No Lagrangian. Discrete steps drift by
  $O(\eta^2)$ relative to the charge, and Adam and momentum break the law.\\
\textbf{Gauge} & The symmetry group of a parametrisation (\thmref{Def.}{I.5}, \bgref{D.21}); for LoRA, $\GL_r$ acting by
  $(Bg^{-1},gA)$. & Generic gauges and fibre dimensions (\thmref{Prop.}{II.12}); the part each optimizer respects,
  $B_r\subset\Ort(r)\subset\GL_r$ (\thmref{Thm}{III.12}); the gauge dependence of factor-wise merging
  (\thmref{Prop.}{V.4}). & No connection or curvature. Fibres are gauge orbits only on the full-rank stratum; away from it
  the fibres jump.\\
\textbf{Moduli} & A set of isomorphism classes (\bgref{C.20}). & Three strata of full-rank LoRA pointings are classified,
  as a set, by $\Gr(r,n)\sqcup\Gr(r,m)\sqcup\MM_r$ (\thmref{Thm}{II.2}). & No topology: the natural quotient topology is
  not a disjoint union ($\mathrm T_A$ lies in the closure of $\mathrm S$). No stack is built, and the strata miss some
  pointings (HF's \texttt{orthogonal} init, defect pointings).\\
\textbf{Fibre} & The preimage $\rho^{-1}(\theta)$, the parameters that give the same weights, as in the title. & Fibre
  dimensions (\thmref{Prop.}{II.12}(a)); which points of the apex fibre $\{BA=0\}$ an optimizer can tell apart
  (\thmref{Cor.}{III.13}). & $\rho$ is no fibre bundle; the fibre dimension jumps between strata.\\
\textbf{Lens} & Backpropagation as the 2-functor $\Para(\mathsf R)$, which sorts tangent-side from cotangent-side methods
  (\bgref{B.17}). & The degree-one collapse: GaLore with weight decay $0$ is one-sided LoRA with a moving frame and
  carried optimizer state (\thmref{Prop.}{IV.2}). & The duality is exact only at degree one, for optimizers that do not
  read the parameter tensor, with no weight decay; decoupled decay, default Adafactor and full-gradient clipping break
  it.\\
\end{longtable}
\endgroup

\paragraph{Credit}
$\Para$ and the lens reading of backpropagation are due to \citet{fong2017backprop} and \citet{cruttwell2021categorical}.
The ETHER control study that tempers the knowledge-retention reading of orthogonal methods is \citet{bini2024ether}.

Three scopes travel with every claim. Expressivity is a different property from learnability. Only
Exposé~\ref{sec:fibres} is dynamical, and each of its results names its optimizer. Mergeability is proved box by box
under explicit side conditions (position-uniformity, no weight sharing, bias slots); the rewrite rules give sufficient
conditions only, and whole-network statements are open.

Most of the mathematics is real algebraic geometry, invariant theory, Lie generation and gradient-flow ODEs. The category
theory does its own work in seven places: the interval [frozen, full FT] and simulation; invariants finer than the
image; a product that yields LoRA-XS, and meets that organise others; base change along maps that change the function;
the lens; naturality of gradient descent and of combination rules; and rank as a min-cut.

Some statements rest on less than a full proof. The equalities in the dimensions of LoHa (\thmref{Proposition}{II.12})
and VeRA are checked numerically, and the bounds are proved; for VeRA with $r\le n$ and generic frozen factors the
equality also follows from the toric-cone description of \bgref{C.14}, and the generic equality for DoRA
(\thmref{Proposition}{II.8}) is proved. The case of unconstrained generators in corollary 4 of \thmref{Theorem}{V.3} is
recorded as a proof sketch, which Appendix~\ref{proof:V.3} completes. The continuous case of
\thmref{Proposition}{V.4}(a) and the converse in \thmref{Corollary}{III.13} are proved in Appendix~\ref{app:proofs}, and
\thmref{Proposition}{IV.3} is argued around Frobenius' theorem (\bgref{D.25}). The ancillary script
\texttt{framework\_checks.py} re-runs the main numbers in about a minute (Appendix~\ref{app:repro}).

% ----------------------------------------------------------------------------------------------------------------------
\subsubsection{Open problems}\label{sec:analogy-open}

Each problem grew out of a proof that stopped short, and each is stated so that a solution would be recognisable.

\begingroup\emergencystretch=2em
\begin{enumerate}[label=\arabic*.,leftmargin=1.8em,itemsep=4pt]
  \item \textbf{LoHa's dimension.} \thmref{Proposition}{II.12} proves
    \begin{equation}\label{eq:open-loha}
      d(\mathrm{LoHa}_{r_1,r_2})\le\min\big(mn,\ r_1(m+n-r_1)+r_2(m+n-r_2)-(m+n-1)\big),
    \end{equation}
    with equality checked only numerically. Prove it in general, and decide whether the Hadamard variety is closed.
  \item \textbf{Initialisations, dynamically.} \thmref{Theorem}{II.2} classifies three strata by expressivity, and
    \thmref{Corollary}{III.13} shows that Adam sees $A_0$ only up to signed row permutations. Classify the apex fibre
    modulo $B_r$ for Adam, stratified by the charge $\Phi_0$ (\thmref{Corollary}{III.6}), including the degenerate
    pointings outside the three strata of \thmref{Theorem}{II.2}.
  \item \textbf{HRA's apex in practice.} HRA's paired initialisation, HF PEFT's default, starts at an apex
    (\thmref{Theorem}{II.7}(b), for injective $W_0$ and even $r\le n-2$), with $\dim S_1=kn-\tfrac{k(k+1)}2$, $k=r/2$,
    against $d=rn-\tfrac{r(r+1)}2$. Decide whether this apex is a practical first-order bottleneck that slows early
    training, by comparison with a regular initialisation at matched budget.
  \item \textbf{Charts for meets.} LoRA-XS (\thmref{Proposition}{I.8}(a)) and HRA's image (\thmref{Theorem}{II.7}(a))
    are meets with a polynomial chart; the meet of LoRA with output scaling (\thmref{Proposition}{I.8}(b)) is a union of
    coordinate planes. Characterise the pairs of named methods whose meet is parametrisable by a polynomial chart.
  \item \textbf{The frontier.} Incomparable pairs are common: split against zero-init LoRA (\thmref{Theorem}{II.2}(d)),
    Kronecker sums or GraLoRA against LoRA (\thmref{Theorem}{II.11}, \thmref{Proposition}{II.12}(c)). Find the size-$k$
    frontier of $\Meth$, whose maximal elements are pairwise incomparable.
  \item \textbf{An Adam analogue of the crossover.} \thmref{Theorem}{III.5} and \thmref{Theorem}{III.9} hold for
    gradient flow. Find a conserved or monotone quantity for sign-like updates, whose symmetry is only $B_r$
    (\thmref{Theorem}{III.12}).
  \item \textbf{Prefixes and prompts.} For a vanilla softmax head with $W_o\ne0$ and $W_q\ne0$, a generic prefix is not
    box-locally mergeable (\thmref{Theorem}{VI.5}(f)); prove or refute its non-mergeability for the whole network. In causal decoders every soft prompt is a deep prefix
    (\thmref{Theorem}{VI.5}(e)). \citet{petrov2023prompting} and \citet{li2021prefix} assert that prompt tuning is
    strictly weaker than deep prefix tuning (P-tuning v2, \citealp{liu2021p}), and the inclusion is strict in a
    one-layer, one-token example. Prove the generic strictness, and settle encoders.
\end{enumerate}
\endgroup

Conjectures \hyperref[res:A]{A}, \hyperref[res:C]{C} and \hyperref[res:D]{D} are open as well. They appear in
Exposé~\ref{sec:predictions} as Predictions \hyperref[res:X.4]{X.4}, \hyperref[res:X.11]{X.11} and
\hyperref[res:X.2]{X.2}.

% ----------------------------------------------------------------------------------------------------------------------
\subsubsection{Debts: known results we re-derive}\label{sec:analogy-known}

Much of what this paper proves was known in another language or setting. A re-derivation earns its place only if the new
language adds something, so Table~\ref{tab:known} states what we add next to each debt.

\begingroup
\footnotesize
\renewcommand{\arraystretch}{1.18}
\setlength{\LTpre}{6pt}\setlength{\LTpost}{8pt}
\begin{longtable}{@{}>{\raggedright\arraybackslash}p{2.35cm}>{\raggedright\arraybackslash}p{6.15cm}>{\raggedright\arraybackslash}p{5.6cm}@{}}
\caption{Known results we re-derive, whom they are due to, and what is added here.}\label{tab:known}\\
\toprule
\headrow \textbf{Here} & \textbf{Due to} & \textbf{Added here}\\
\midrule
\endfirsthead
\multicolumn{3}{@{}l}{\footnotesize\itshape Table~\ref{tab:known}, continued}\\
\toprule
\headrow \textbf{Here} & \textbf{Due to} & \textbf{Added here}\\
\midrule
\endhead
\midrule
\multicolumn{3}{@{}r}{\footnotesize\itshape continued on the next page}\\
\endfoot
\bottomrule
\endlastfoot
\thmref{Def.}{I.1}, Exposé~\ref{sec:lens} & \citet{fong2017backprop}; \citet{cruttwell2021categorical} & The pointed
  slice of methods; tangent and cotangent sides.\\
\thmref{Thm}{II.1} & The PiSSA-to-LoRA conversion of \citet{meng2024pissa}, in HF PEFT \citep{xmangrulkar2022peft} & The
  rank-$2r$ bound for every exact pointing, and from any start.\\
\thmref{Thm}{II.5} & Classical: the first fundamental theorem for $\Ort(n)$; Witt's extension theorem (\bgref{D.16}) &
  Which invariant each group-type method keeps.\\
\thmref{Thm}{II.10} & The standard cut-rank bound for tensor networks (\bgref{B.8}) & A rank census of PEFT plugs.\\
\thmref{Thm}{II.11} & \citet{loan1993approximation}; the operator-Schmidt decomposition (\bgref{B.12}) & Kronecker sums
  as LoRA across another cut; the LoKr inclusion criterion.\\
\thmref{Prop.}{II.12}, TT/Tucker row & \citet{xkoch2010dynamical}; \citet{xuschmajew2013geometry};
  \citet{xholtz2012manifolds} & The same count, at minimal ranks, beside other PEFT families.\\
\thmref{Prop.}{III.4}, \thmref{Thm}{III.12}(c) & \citet{xtong2021accelerating}; \citet{zhang2024riemannian}; LoRA-Pro
  \citep{wang2024lorac}; LoRA-FA \citep{zhang2023lora}; descent to the quotient: \citet{mishra2012fixed},
  \citet{xmishra2016riemannian} & One cometric table for (S)GD; the hierarchy $B_r\subset\Ort(r)\subset\GL_r$.\\
\thmref{Thm}{III.5} & \citet{zhao2022symmetries}; \citet{marcotte2023abide}; the principle: \citet{du2018algorithmic},
  \citet{arora2018optimization}, \citet{kunin2020neural} & The formally immediate remark that the signal need not be a
  gradient; PEFT instances (DoRA's detached norm, clipping, VeRA, AdaLoRA without its regulariser, HRA).\\
\thmref{Prop.}{III.7} & \citet{xkempf1979length}, in the real form of \citet{xrichardson1990minimum};
  \citet{xsrebro2004maximum} & Balance as the Kempf--Ness slice of LoRA's gauge; two decay conventions, two slices
  (\thmref{Cor.}{III.6}).\\
\thmref{Thm}{III.9} & \citet{xtarmoun2021understanding}; \citet{xmin2021explicit}; \citet{xdomine2025lazy};
  \citet{xkunin2024get} & The PEFT reading of $\sigma^*$, its shift under LoRA+, PiSSA's balance, the Kaiming band.\\
\thmref{Thm}{III.10}, \thmref{Cor.}{III.11} & \citet{schulman2025lora}, for LoRA+ \citep{hayou2024lora} & Any
  multihomogeneous method; the SGD $c^2$ law; per-block $\varepsilon$, decay and clipping.\\
\thmref{Prop.}{IV.2} & \citet{xtorrobahennigen2025duality}; for random projections, Flora \citep{hao2024flora} & The
  projector-by-schedule table of cotangent methods; optimizer state carried across re-basings.\\
\thmref{Prop.}{IV.3} & Frobenius' theorem (\bgref{D.25}) & When a $\theta$-dependent anchor foliates weight space by
  images of forward methods; GaLore, LISA and MeZO keep a fixed anchor per period, so the test does not separate
  them.\\
\thmref{Lemma}{VI.2} & \citet{xsussmann1992uniqueness}; \citet{xchen1993geometry}; \citet{xgodfrey2022symmetries} &
  Which symmetries a method may ignore; the monomial group of gated FFNs.\\
\thmref{Thm}{VI.3} & (IA)$^3$ merging: \citet{liu2022few}, Sec.~3.3; SSF's re-parameterisation \citep{lian2022scaling};
  LoRA's linearity \citep{hu2021lora}; scale migration through SwiGLU: AWQ \citep{xlin2024awq} & One calculus of
  sufficient conditions, with explicit side conditions.\\
\thmref{Thm}{VI.5} & The identity is standard; (b) \citet{he2021unified}; (c) and the prompt-to-prefix map (e)
  \citet{petrov2023prompting}, Sec.~2.2, who also show that prefix-tuning a pretrained transformer can be a universal
  approximator \citep{petrov2024prompting} & The homomorphism-then-perspective organisation; explicit position hypotheses
  and the slice typing of (e).\\
\end{longtable}
\endgroup

\paragraph{Credit}
We also build on work without re-deriving it. \citet{malladi2022kernel} found the lazy kernel of random-$A_0$ LoRA close
to that of full fine-tuning, which \thmref{Observation}{I.6} reads through $S_1$. The Init[B] scheme in the table of
\thmref{Theorem}{II.2} is from \citet{hayou2024impact}. With a fixed support, the meet of \thmref{Proposition}{I.8}(b)
is (IA)$^3$ masked to $r$ neurons; compare the sparse activation scaling of \citet{xstoehr2024activation}.

% ----------------------------------------------------------------------------------------------------------------------
\subsubsection{The referee record}\label{sec:analogy-referee}

The framework was synthesised from four proposals and three referee reports. Every numbered result then went to two
independent referees, who checked it against its proof, the literature and the Hugging Face PEFT source, that is, for
its proof and for its fit to the methods as published and as coded. Of the \nPropositions{} results, 14 stood as first
written, with wording fixes, and 27 were repaired by an added hypothesis, a narrowed claim, a replaced proof or a
withdrawn reading. A second review re-examined 26 of the repaired statements in 48 reports. It raised a major objection
to 18 of them, mostly to a side clause: a missing hypothesis, a remark that went past its proof, or a description of a
method that matched only one implementation. For example, HF PEFT's LoftQ initialiser turned out to ignore the LoRA
scale $\alpha/r$ (\thmref{Proposition}{V.2}). In all it found a false or overreaching clause in 19 statements, and those
were narrowed again. Only the final statements appear in this paper; the referee notes of both rounds are in the
ancillary file \texttt{propositions.json}. Table~\ref{tab:referee} records what changed, result by result.

\begingroup
\footnotesize
\renewcommand{\arraystretch}{1.15}
\setlength{\LTpre}{6pt}\setlength{\LTpost}{8pt}
\begin{longtable}{@{}>{\raggedright\arraybackslash}p{3.45cm}>{\raggedright\arraybackslash}p{1.75cm}>{\raggedright\arraybackslash}p{8.9cm}@{}}
\caption{The referee record: the status of each numbered result after both rounds of review, and what the referees
changed.}\label{tab:referee}\\
\toprule
\headrow \textbf{Result} & \textbf{Status} & \textbf{What the referees changed}\\
\midrule
\endfirsthead
\multicolumn{3}{@{}l}{\footnotesize\itshape Table~\ref{tab:referee}, continued}\\
\toprule
\headrow \textbf{Result} & \textbf{Status} & \textbf{What the referees changed}\\
\midrule
\endhead
\midrule
\multicolumn{3}{@{}r}{\footnotesize\itshape continued on the next page}\\
\endfoot
\bottomrule
\endlastfoot
\thmref{Obs.}{I.4}, universal properties & repaired & Products need a fibre product in $\Diff$ and can fail to exist
  ($\mathrm{LoRA}_1\times\text{Frozen}$ on $1\times1$ weights).\\
\thmref{Obs.}{I.6}, $S_1$ is functorial & verified & Monotonicity stated only along pointed morphisms.\\
\thmref{Obs.}{I.7}, additive monoidal structure & repaired twice & $r\mapsto\mathrm{LoRA}_r$ is strong monoidal only
  unpointed or at degenerate pointings ($\dim S_1$ is $mr$ against $2mr$); HF's \texttt{cat} merge described exactly.\\
\thmref{Prop.}{I.8}, a product and a meet & verified & LoRA-XS needs $\rk W_0\ge r$; the set-level meet in (b) is no
  new method.\\
\thmref{Thm}{II.1}, pointings of an additive method & verified & Stated for exact pointings, with a robust form;
  credited to \citet{meng2024pissa}.\\
\thmref{Thm}{II.2}, three strata of LoRA initialisations & repaired & Full-rank pointings only, and only as a set; HF's
  \texttt{orthogonal} init and QLoRA/LoftQ defects lie outside.\\
\thmref{Prop.}{II.3}, first-order rank formula & repaired & Scale and gradient flow made explicit; Adam's first step has
  an $\ell_1$ rate; LoRA-GA's $A_0$, not its $B_0$, maximises the rate.\\
\thmref{Lemma}{II.4}, zero gate & verified & Hypotheses stated; a function-level variant for LLaMA-Adapter; AdaLoRA
  instance corrected.\\
\thmref{Thm}{II.5}, complete invariants & verified & Methods that can create zeros ((IA)$^3$, HiRA) keep only ``zeros of
  $W_0$ stay zero''.\\
\thmref{Cor.}{II.6}, what orthogonal methods never do & repaired twice & Exactly orthogonal input-side maps only; BOFT
  keeps cosines up to sign; HF's Cayley--Neumann OFT and GOFT's scaler excluded.\\
\thmref{Thm}{II.7}, HRA as a meet; its apex & repaired twice & An image-level meet with $W_0\cdot\SO(n)$, no product;
  strictly above LoRA $\cap$ OFT only for exactly orthogonal OFT.\\
\thmref{Prop.}{II.8}, DoRA's image & repaired twice & $\mathrm{Diag}_m\cdot(W_0+\Mr)$ for $W_0$ without zero rows; the
  rank formula needs nonzero magnitudes; Conjecture~D made falsifiable.\\
\thmref{Prop.}{II.9}, the torus ladder & repaired & $\mathrm{RoAd}_1$ is (IA)$^3$ over $\mathbb C$; only zeros of $W_0$
  persist; (IA)$^3$'s FFN vector acts on the input of $W_{\mathrm{down}}$.\\
\thmref{Thm}{II.10}, cut-rank bound & verified & A large min-cut is necessary for high rank and is not sufficient.\\
\thmref{Thm}{II.11}, Kronecker incomparability & repaired twice & $\kappa_1$ is the maximal Kronecker rank for aligned
  splits; LoKr with $r'\le\min(m_2,n_2)$ lies in $\mathrm{LoRA}_r$ iff $r\ge\min(m_1,n_1)\,r'$.\\
\thmref{Prop.}{II.12}, gauge census & repaired twice & GraLoRA's maximum rank is $\min(kr,m,n)$; LoHa loses dimensions
  to $\mathrm{LoRA}_{2r}$ only when $2r^2<m+n-1$.\\
\thmref{Prop.}{II.13}, transport of structure & repaired twice & Any $T\in\GL(\Theta)$; commutes with rebasing for
  additive families and fixed $T$; FourierFT needs no conjugate pair.\\
\thmref{Obs.}{II.15}, descent & repaired twice & Reachable sets only, for box-wise methods without shared frames;
  residual rotations need a gain-fused pre-norm model.\\
\thmref{Prop.}{II.16}, covariant splits & repaired twice & MiLoRA and CorDA-KPM are bottom-$r$ splits; trace-relative
  damping keeps $\CO(n)$; shipped CorDA code is not covered.\\
\thmref{Obs.}{III.2}, chain rule for cometrics & verified & Scope stated: gradient flow, no LoRA dropout, constant
  rates.\\
\thmref{Thm}{III.3}, naturality of gradient descent & verified & The title ``natural only along isometries'' was false in
  general and was replaced.\\
\thmref{Prop.}{III.4}, cometrics that descend & repaired & Undamped (S)GD only; the published methods use AdamW or
  damping.\\
\thmref{Thm}{III.5}, Noether charges & repaired twice & Isometries need the signal to factor through $(t,\rho(q))$;
  relabelled known; the per-channel-rate remark needs distinct rate products.\\
\thmref{Cor.}{III.6}, charge at initialisation; decay & repaired & The decay commentary was reversed: coupled decay tends
  to $\mu=0$, per-time decay to $\Phi=0$.\\
\thmref{Prop.}{III.7}, Kempf--Ness slice & verified & Uniqueness up to $\Ort(r)$ proved; the Loewner floor renamed.\\
\thmref{Prop.}{III.8}, balance--escape trilemma & verified & Balance generalised to $B_0\T B_0=cA_0A_0\T$.\\
\thmref{Thm}{III.9}, isotropic-charge dynamics & verified & Pointwise on the rank-$r$ locus; the Kaiming crossover is a
  band.\\
\thmref{Thm}{III.10}, the hyperparameter torus & verified & Hypotheses explicit; extended to schedules, per-block
  $\varepsilon$, decay and clipping.\\
\thmref{Cor.}{III.11}, LoRA+ and rsLoRA & verified & A shared warm-up keeps the torus; linear modules only.\\
\thmref{Thm}{III.12}, optimizer gauge hierarchy & repaired twice & Adam respects $B_r$, positive diagonals only with
  per-channel rates, $\varepsilon$ and decay; damping cuts $\GL_r$ to $\Ort(r)$.\\
\thmref{Cor.}{III.13}, what zero-$B$ chooses & repaired & The apex fibre is $\{BA=0\}$; the $\GL_r$ case needs
  equivariance along the whole trajectory.\\
\thmref{Prop.}{IV.2}, degree-one collapse & repaired & Decoupled decay, default Adafactor and full clipping break it;
  GaLore carries optimizer state.\\
\thmref{Prop.}{IV.3}, integrability & repaired twice & The ``only if'' was false; the top-$r$ gradient field is
  non-involutive for every $r\ge1$.\\
\thmref{Obs.}{IV.4}, the backprop cone & repaired twice & An upper bound; forward mode pays off only for $k=O(1)$.\\
\thmref{Prop.}{V.2}, quantisation and splitting & repaired twice & Three orders, three defects; HF's LoftQ init ignores
  the LoRA scale; QA-LoRA merges on-grid.\\
\thmref{Thm}{V.3}, rebasing closure & repaired twice & BOFT reaches $\SO(n)$ only at full depth; HF's Cayley--Neumann
  merges contract.\\
\thmref{Prop.}{V.4}, combination rules & repaired twice & The ``iff'' needs full rank or continuity; factor-wise DARE is
  gauge-dependent too.\\
\thmref{Lemma}{VI.2}, activations fix the basis & verified & Gated FFNs have the full monomial group.\\
\thmref{Thm}{VI.3}, the merge calculus & repaired twice & Sufficient conditions only; RoPE needs scales constant on
  rotary pairs.\\
\thmref{Prop.}{VI.4}, linear insertions, ReFT & repaired twice & LoReFT has dimension $r(2d+1-r)$ and a pyreft pointing
  defect; the serial-adapter criterion needs $\rk A_0=r$.\\
\thmref{Thm}{VI.5}, attention as a homomorphism & repaired twice & Rigidity is per layer; the prompt-to-prefix map
  credited to \citet{petrov2023prompting}.\\
\end{longtable}
\endgroup

A repair is a record of an error that was caught. We keep the record public so that a reader can weigh each result by
its history, and so that the next error, when someone finds it, has an obvious place to go.
